% !TEX root = ../../tree_notes.tex
% =====================================================================
%  8.1  Why Balance? The Rotation                sec:c06-rotations
%  Exports: def:c06-heightcost  thm:c06-rotation  def:c06-rotcases
%           fig:c06-rot  lst:c06-rotcode  ex:c06-rot-run
%           prob:c06-treecons  lst:c06-treecons-soln   (practice: LC 1382)
% =====================================================================
\section{Why Balance? The Rotation}\label{sec:c06-rotations}

\begin{previouslybox}
\begin{enumerate}[nosep, leftmargin=1.6em]
  \item What is the BST invariant, and why does one comparison \emph{discard a whole subtree} during
        a search? (\chref{c05}, \defref{c05-prop}, \propref{c05-search})
  \item Every Chapter-7 routine -- search, insert, the three-case delete -- costs how much, written
        in which single parameter? (\bigO{h}; \secref{c05-insert-delete})
  \item Insert $1,2,\dots,n$ into a plain BST. What shape do you get, and what does $h$ become?
        (a chain, $h=n-1$; \secref{c05-contests}' warning, \probref{c05-summit}' cliffhanger)
  \item Which traversal of a BST lists the keys in sorted order, and why will we lean on it being
        \emph{invariant under reshaping}? (\chref{c03}, \thmref{c05-inorder})
\end{enumerate}
\end{previouslybox}

\begin{topicintro}
Chapter~7 proved a beautiful thing and then quietly made you pay for it: a BST answers
membership, insertion, deletion and every ordered query in $\bigO{h}$, \emph{where $h$ is the height}
-- and $h$ is a property of the \emph{input order}, not of you. Insert sorted keys and the tree is a
linked list with spare pointers; the Summit \probref{c05-summit} went quadratic for exactly this
reason. This section installs the fix. There is only \emph{one} legal local operation on a binary
search tree -- the \term{rotation} -- and its magic is that it changes the shape (and hence the
height) while leaving the sorted inorder line \emph{byte for byte identical}. Once you accept the
rotation as the atomic move, every balanced tree in the rest of the chapter is a \emph{policy}: a
rule for when to rotate and how hard to keep things short. We finish on Codeforces 675D, which hands
you a degenerate insertion order and asks for each node's parent -- the cleanest possible way to feel
both the $\bigO{n^2}$ trap and the $\bigO{n\log n}$ escape.
\textbf{Prerequisites.} All of Chapter~7 (\chref{c05}); the descent, the invariant, and the
$\bigO{h}$-cost warning are assumed. Inorder from Chapter~3 (\chref{c03}) is used as a black box,
never re-derived.

\textbf{After this section you can}
\begin{itemize}[nosep]
  \item state precisely why every BST operation is $\bigO{h}$ and why $h=\bigO{\log n}$ is the goal; \hfill\lv{2}
  \item perform a left and a right rotation and prove the inorder sequence is unchanged; \hfill\lv{2}--\lv{3}
  \item classify a height violation into LL / RR / LR / RL and pick the one- or two-rotation fix; \hfill\lv{3}
  \item solve BST-construction-under-degenerate-order problems with a set, escaping $\bigO{n^2}$. \hfill\lv{4}
\end{itemize}
\end{topicintro}

\begin{recognitionbox}[Pattern card -- the shape is the problem]
\reckey{a search tree's \emph{efficiency}, not its answer, depends on insertion order; or you are
asked to \emph{rebuild/reshape} a BST while preserving its sorted contents.}
\recuse{
  \item ``insert these keys, tree must stay fast'' / keys arrive sorted -- balance by construction
        (the whole chapter; start \secref{c06-avl});
  \item ``give each node's parent in a BST built by this order'' -- CF 675D, the set-neighbour trick
        (\probref{c06-treecons});
  \item ``turn a skewed BST into a balanced one'' -- rebuild from inorder (Practice, LC 1382);
  \item any rotation-tracing or ``apply one left rotation at $x$'' warm-up -- \defref{c06-rotcases}.
}
\recnotuse{
  \item the tree is \emph{already} balanced and you just query it -- that is plain Chapter~7 descent
        (\chref{c05});
  \item keys are dense integers and you only need prefix ranks -- Fenwick \chref{c14} skips trees
        entirely;
  \item you need whole-subtree link/cut, not single-node rebalancing -- splay/link-cut, \chref{c23}.
}
One-line test: \emph{is the \textbf{cost} of a correct answer the thing at stake -- a height you must
force down without disturbing the sorted order? Then you need rotations, and this chapter.}
\end{recognitionbox}

\subsection{Intuition}

Picture the running chain $C$ of \figref{c06-chain}: keys $1,2,\dots,7$ inserted in increasing order.
Every new key is larger than everything present, so it always goes right -- a straight diagonal of
seven nodes, height $6$. A search for $7$ now touches all seven nodes: the tree has \emph{forgotten}
how to binary-search and degenerated into a linked list. Yet the \emph{contents} are perfectly fine --
read the chain's inorder (left, node, right) and you still get $1,2,3,4,5,6,7$, sorted. The problem is
purely geometric: same data, bad shape.

So we want a move that is to a BST what \emph{reshuffling a bookshelf} is to a library: change which
shelf is the central support without re-ordering a single book. That move is the \term{rotation}. It
re-parents \emph{three} pointers around one node, so its cost is $\bigO{1}$, and -- crucially --
\thmref{c06-rotation} guarantees it cannot disturb the inorder line. Do enough of them and the chain
$C$ becomes the balanced twin on the right of \figref{c06-chain}: height drops from $6$ to $2$, every
descent is $\bigO{\log n}$ again, and \emph{not one key moved in the sorted order}.

\begin{figure}[H]
\centering
\begin{tikzpicture}[diagram, scale=0.95, transform shape]
  % left rotation about x:  before
  \node[tnode] (xa) at (0,0) {$x$};
  \node[tleaf] (aA) at (-1.2,-1.3) {$A$};
  \node[tnode] (ay) at (1.2,-1.3) {$y$};
  \node[tleaf] (aB) at (0.3,-2.6) {$B$};
  \node[tleaf] (aC) at (2.1,-2.6) {$C$};
  \draw (xa) -- (aA); \draw (xa) -- (ay); \draw (ay) -- (aB); \draw (ay) -- (aC);
  \node[annot, font=\scriptsize] at (0.5,1.0) {left rotate about $x$};
  \node[annot, font=\scriptsize\bfseries] at (0,-4.0) {before};
  % arrow
  \draw[flow, double] (3.4,-1.4) -- (5.0,-1.4);
  % after
  \node[tnode] (yb) at (6.6,0) {$y$};
  \node[tnode] (bx) at (5.0,-1.3) {$x$};
  \node[tleaf] (cC) at (8.0,-1.3) {$C$};
  \node[tleaf] (bA) at (4.0,-2.6) {$A$};
  \node[tleaf] (bB) at (6.0,-2.6) {$B$};
  \draw (yb) -- (bx); \draw (yb) -- (cC); \draw (bx) -- (bA); \draw (bx) -- (bB);
  \node[annot, font=\scriptsize\bfseries] at (6.5,-4.0) {after};
  \node[annot, align=left, font=\scriptsize] at (8.6,-2.4)
    {inorder both:\\ $A,x,B,y,C$\\ \emph{identical} --\\ only $x,y,B$ re-linked};
\end{tikzpicture}
\caption{A left rotation about $x$ (right child $y$): $y$ is promoted to $x$'s former spot, $x$
becomes $y$'s left child, and $y$'s old left subtree $B$ is handed down to become $x$'s right child.
Exactly three subtree pointers move; the inorder sequence $A,x,B,y,C$ is preserved verbatim, which is
\thmref{c06-rotation}. A right rotation is the exact mirror image.}\label{fig:c06-rot}
\end{figure}

\subsection{Definitions and Notation}

\begin{definition}{Height is the whole cost}{c06-heightcost}
For a BST of $n$ distinct keys and height $h$, every one-descent operation of Chapter~7 --
membership, the failed-search insert, the successor-splice delete, a kth-smallest walk -- costs
$\bigO{h}$. A tree is \term{height-balanced} if there is a constant $c$ with $h \le c\,\lg(n+1)$
regardless of insertion order. Because $2 \le h+1 \le n$ always, ``balanced'' means the descent is
guaranteed logarithmic, and \emph{un}balanced means it can be linear: the entire difference between
the Summit's $\bigO{n^2}$ failure and a $\bigO{n\log n}$ success.
\end{definition}

\noindent\textbf{In plain words.} Stop thinking ``$\bigO{n}$'' for a search. Think ``$\bigO{h}$'',
and make keeping $h$ tiny the \emph{data structure's job}, so the algorithm never has to.

\begin{theorem}{A rotation preserves the inorder sequence}{c06-rotation}
Let $x$ be a node with right child $y$, left subtree $A$, and $y$'s left subtree $B$, right subtree
$C$ (so the subtree's inorder is $A,x,B,y,C$). The left rotation about $x$ producing $y\to(x\to(A,B),C)$
yields the \emph{same} inorder sequence $A,x,B,y,C$, and its inverse (right rotation about $y$) restores
the original. Symmetrically for a right rotation. Hence rotations preserve both the \emph{sorted
order} and the BST property of the affected subtree.
\end{theorem}
\begin{proofbox}
Read the inorder of the \emph{after} picture. The new root is $y$; its left subtree is rooted at $x$,
whose own children are $A$ (left) and $B$ (right), and whose right child under $y$ is $C$. Visit $y$'s
left subtree first: within $x$, visit $A$, then $x$, then $B$ -- giving $A,x,B$. Then visit $y$. Then
visit $y$'s right subtree $C$. Total: $A,x,B,y,C$ -- term for term the \emph{before} inorder. No key
crossed any other, so the BST order (every subtree a contiguous sorted run) is intact on both sides.
Undoing is the mirror right rotation about $y$, which we just showed preserves inorder too, so the two
operations are exact inverses. \end{proofbox}
\begin{insightbox}
\thmref{c06-rotation} is the \emph{licence to reshape}. It says geometry and semantics are decoupled
in a BST: the sorted line is fixed, and the \emph{tree that realises it} is a choice. Height is a
property of the realisation, not the data -- so we are free to re-choose it (rotate) whenever we
like, and every balancing scheme in Sections 8.2--8.5 is exactly a rule for when to spend a rotation.
\end{insightbox}

\begin{definition}{The four violation cases: LL, RR, LR, RL}{c06-rotcases}
A node $v$ is \term{unbalanced} when its two subtrees differ in height too much (the precise threshold
defines AVL vs red-black vs nothing; here ``one side is a chain and the other short'' suffices). Look
at the heavy child and which side \emph{it} leans:
\begin{itemize}[nosep]
  \item \textbf{LL} -- heavy child leans Left: one \emph{right} rotation about $v$ fixes it.
  \item \textbf{RR} -- heavy child leans Right: one \emph{left} rotation about $v$.
  \item \textbf{LR} -- heavy child leans Right (zig-zag): first a left rotation about the child turns
        it into LL, then a right rotation about $v$. (This is $Z$, \figref{c06-zigzag}.)
  \item \textbf{RL} -- heavy child leans Left (zig-zag): first right about the child, then left about $v$.
\end{itemize}
The single-rotation cases are \emph{zig-zig} (both bends the same way); the double are \emph{zig-zag}
(bends differ), needing exactly two rotations and three relinkings in total.
\end{definition}

\noindent\textbf{In plain words.} A same-direction lean is one rotation; an opposite-direction lean
(the child sticks out the far side) is \emph{two}, because the first rotation converts the zig-zag
into a zig-zig, then the second finishes. Node $Z$ needs two; you cannot fix it with one, which is
the whole lesson of \figref{c06-zigzag}.

\subsection{The algorithm: one rotation, written once}

Every balancing scheme reuses the same two four-line routines. Keep pointers to the three moving
subtrees ($A,B,C$), relink twice, and recompute the two affected nodes' heights -- done. That
$\bigO{1}$ cost is why a single rotation is ``free,'' and why the schemes differ only in
\emph{decision}, never in \emph{mechanics}.

\begin{lstlisting}[caption={Left and right rotations about a node -- the atomic move reused by AVL, red-black and splay}, label={lst:c06-rotcode}]
// Left rotation about x (x must have a right child y). Returns the new subtree root y.
Node* rotateLeft(Node* x) {
    Node* y = x->right;          // y is promoted
    x->right = y->left;          // B moves under x as its right child
    y->left  = x;                // x becomes y's left child
    updHeight(x); updHeight(y);  // x changed first, then its new parent y
    return y;                    // y is the root of the reshaped subtree
}
// Right rotation about y (mirror image; y must have a left child x). Returns new root x.
Node* rotateRight(Node* y) {
    Node* x = y->left;
    y->left  = x->right;         // B moves under y as its left child
    x->right = y;
    updHeight(y); updHeight(x);
    return x;
}
\end{lstlisting}
\begin{complexitybox}
Time $\bigO{1}$ -- a fixed number of pointer assignments and two height recomputes; memory $\bigO{1}$.
This is the point: rotations are constant-time, so a balancing policy that performs a \emph{constant
number} of them per edit (AVL: $\le 2$; red-black: $\le 2$ rotations plus recolouring; splay: $O(1)$
amortised) keeps every Chapter-7 operation at its $\bigO{\log n}$ ideal. Verdict: the atomic
$\bigO{1}$ move \emph{toward} balance.
\end{complexitybox}

\subsection{Worked Example}

\begin{example}{Rotating the chain $C$ and resolving the zig-zag $Z$}{c06-rot-run}
\emph{Chain $C$, top triple.} Focus on the head $1\to 2\to 3$ (each going right) of $1,2,\dots,7$.
Apply \texttt{rotateLeft(1)}: $2$ is promoted, $1$ becomes $2$'s left child, and $2$'s (empty) left
subtree $B$ becomes $1$'s right. Local shape $1\to 2\to(\text{nothing},3)$ becomes $2\to(1,3)$ -- a
height-$1$ balanced triple instead of a height-$2$ chain. Repeat on the next triple and, after
$\bigO{1}$ rotations per level, the diagonal flattens toward the balanced twin of
\figref{c06-chain}. Inorder is untouched: at every step the run reads $1,2,3,\dots$.\\
\emph{Zig-zag $Z$.} At node $30$ the left subtree ($10\to 20\to 25$, all going right) has height $2$
while the right ($40$) has height $0$: an LL-vs-RL mismatch -- specifically the heavy left child $10$
leans \emph{right}, the \textbf{LR} case. One rotation cannot fix it. First \texttt{rotateLeft(10)}:
$20$ is promoted, giving $30\to(20\to(10,25),40)$ -- now the heavy child $20$ leans \emph{left} at
$30$'s subtree, a clean LL. Then \texttt{rotateRight(30)}: $20$ reaches the top, final shape
$20\to(10, 30\to(25,40))$, height $2$, balanced. Inorder $10,20,25,30,40$ before and after --
\thmref{c06-rotation} honoured twice.
\end{example}
\begin{tracebox}[The two rotations on $Z$, pointer by pointer]
\textbf{Start:} $30$'s left $=10$, $10$'s right $=20$, $20$'s right $=25$, $30$'s right $=40$.
\textbf{rotateLeft(10):} $B := 20$'s left (null). Set $10\to$right $=B=$ null; $20\to$left $=10$.
Now $30\to$left $=20$, $20\to$left $=10$, $20\to$right $=25$. Heavy child of $30$ is now a
\emph{left-leaning} node ($20$ with left child $10$) -- the LL case.
\textbf{rotateRight(30):} $B := 20$'s right ($=25$). Set $30\to$left $=B=25$; $20\to$right $=30$.
Now root $=20$, $20\to$left $=10$, $20\to$right $=30$, $30\to$left $=25$, $30\to$right $=40$. Done:
height $2$, both subtrees of $20$ height $1$. \checkmark Two rotations, inorder $10,20,25,30,40$
preserved -- see \figref{c06-zigzag}.
\end{tracebox}

\subsection{Variations and Twists}

\begin{itemize}[nosep]
  \item \textbf{Double rotation as two singles}: LR $=$ left-then-right and RL $=$ right-then-left
        factor into two single rotations about adjacent nodes -- implement \code{rotateLeft} /
        \code{rotateRight} once and compose; never hand-code a ``double.'' \hfill\lv{3}
  \item \textbf{Rebuild instead of rotate}: dump the BST's inorder into an array and build a balanced
        tree by recursive halving -- the LC 1382 approach; $\bigO{n}$ with no rotations at all.
        \hfill\lv{3}
  \item \textbf{Rotation distance}: any two BSTs on the same keys are connected by $\bigO{n}$
        rotations (Catalan-flavoured; a research-grade fact, cited not used here). \hfill\lv{5}
  \item \textbf{Rotations preserve subtree \emph{sets}}: the three subtrees $A,B,C$ keep their exact
        key contents; only the parentage changes -- this is why an augmented field like
        $\mathit{size}()$ needs \emph{recomputing} on the two moved nodes, not a rebuild. \hfill\lv{4}
\end{itemize}

\subsection{Common Pitfalls}\label{sec:c06-rotations-pitfalls}

\begin{warningbox}[Trap 1 -- rotating the wrong way relative to the lean]
A heavy-left subtree (LL) needs a \emph{right} rotation; a common slip is to left-rotate and deepen
the chain further. Decide the direction from the side that is \emph{too tall}: pull the tall side up.
If after a rotation the height went the wrong way, you rotated about the wrong node or side.
\end{warningbox}
\begin{lstlisting}[caption={\textcolor{warningred}{WRONG}: left-rotate an LL (heavy-left) node -- deepens the left chain}, language=C++, basicstyle=\ttfamily\small]
// v's LEFT subtree is too tall (LL). rotateLeft(v) requires a right child and
// makes the LEFT side even heavier -- the opposite of the fix you need.
Node* badFix(Node* v) { return rotateLeft(v); }   // WRONG direction for LL
\end{lstlisting}
\begin{lstlisting}[caption={\textcolor{successgreen}{RIGHT}: LL is repaired by one right rotation about the anchor}, language=C++, basicstyle=\ttfamily\small]
Node* fixLL(Node* v) { return rotateRight(v); }     // pull the heavy LEFT child up
// LR (zig-zag): first straighten the child, then rotate the anchor.
Node* fixLR(Node* v) { v->left = rotateLeft(v->left); return rotateRight(v); }
\end{lstlisting}

\begin{warningbox}[Trap 2 -- updating heights in the wrong order]
After \texttt{rotateLeft(x)} the node $x$ ends up \emph{below} $y$. If you recompute $y$'s height
before $x$'s, $y$ reads a stale child height. Update the node that moved \emph{down} first, then the
one promoted: exactly the order in \code{rotateLeft} (\coderef{c06-rotcode}). A stale balance
factor silently turns AVL into the same old chain.
\end{warningbox}
\begin{lstlisting}[caption={\textcolor{warningred}{WRONG}: promoted node recomputed before its new child}, language=C++, basicstyle=\ttfamily\small]
x->right = y->left; y->left = x;
updHeight(y);            // y still sees x's OLD height here -- WRONG
updHeight(x);
\end{lstlisting}
\begin{lstlisting}[caption={\textcolor{successgreen}{RIGHT}: fix the lower node $x$ first, then its parent $y$}, language=C++, basicstyle=\ttfamily\small]
x->right = y->left; y->left = x;
updHeight(x);            // x's children are final
updHeight(y);            // now reads x's corrected height -- RIGHT
\end{lstlisting}

\begin{warningbox}[Trap 3 -- simulating degenerate insertion with an $\bigO{h}$ descent]
``Insert $n$ keys into a BST and report each parent'' done by walking down the tree is
$\bigO{h}$ per key, i.e.\ $\bigO{n^2}$ on the chain $C$ -- and CF 675D's $n\le 10^{5}$ will
\tle. The escape (\probref{c06-treecons}) needs \emph{no tree at all}: a \code{std::set} finds each
key's neighbours in $\bigO{\log n}$ and the parent falls out of their insertion times.
\end{warningbox}

\subsection{Problem Lab}

The section's anchor is Codeforces 675D -- build a BST by a given insertion order and report each
node's parent. Simulated honestly it is the degenerate-height trap in the raw ($\bigO{n^2}$); solved
cleverly it is a two-line set query. It is chosen precisely because it \emph{is} the reason this
chapter exists.

% ---------------------------------------------------------------------------
\solvedproblem{c06-treecons}{Tree Construction}{Codeforces 675D}
  {https://codeforces.com/problemset/problem/675/D}{CF 1700 · L4}
  {BST insertion order, set-based neighbour lookup, degenerate-height escape, later-inserted parent}

\begin{problemstatement}[Tree Construction]
You are given a sequence $a_1,\dots,a_n$ of $n$ \textbf{distinct} integers, the order in which they
are inserted into an initially empty binary search tree (each inserted by the usual
$\bigO{h}$ descent of Chapter~7). For every element \emph{except the first}, output the value of its
\textbf{parent} node in the finished BST. $3\le n\le 10^{5}$ and $1\le a_i\le 10^{9}$. Output the
parents one per line, in the order the corresponding children were inserted.
\end{problemstatement}

\begin{attemptgate}[Attempt it yourself first -- \lv{4}, give yourself 25 minutes]
The honest simulation (actually build the tree) is correct but $\bigO{n^2}$ on sorted input -- the
chain $C$ again. The insight is to \emph{not} build the tree.
\peek{When you insert a fresh key $x$ into a BST, its parent is an \emph{already-present} key. Which
two candidates could it possibly be, in terms of the sorted order of the keys already inserted?}
\nudge{$x$'s parent must be its predecessor or its successor \emph{among the keys already present}
(the \code{std::set} can find both in $\bigO{\log n}$). The one inserted \emph{later} is deeper and
becomes the parent -- the earlier one already had the later on the path.}
\fullhint{Keep a \code{set} of present keys and a \code{map} key$\to$insertion-index \code{when}. For
$x$: \code{auto r = s.lower\_bound(x)}; the successor is \code{*r} (when \code{r} is not \code{s.end()}), the predecessor is
\code{*prev(r)} (when \code{r} is not \code{s.begin()}). Parent $=$ whichever of them has the \emph{larger} \code{when}.
Print it for $i\ge 2$, then \code{s.insert(x)} and set \code{when[x]=i}. Every step is
$\bigO{\log n}$, total $\bigO{n\log n}$ -- immune to the degenerate order that sinks the simulation.}
\end{attemptgate}

\solstep{Understanding the Problem}
We must report parent pointers of a BST defined purely by insertion order. The naive route walks the
real tree down from the root -- but on sorted data that tree is the chain $C$, so each walk is
$\bigO{n}$ and the whole is $\bigO{n^2}$, too slow for $n=10^{5}$. The reframing: the finished shape
is a \emph{function of the sorted contents and the arrival times}, not of any pointer we must build.
A new key $x$ hangs as a leaf; its parent is the last-inserted of the two keys that $x$ sits
\emph{between} in sorted order. So we only need, for each $x$, its predecessor and successor among the
already-inserted keys and their insertion order -- a job for \code{std::set} in $\bigO{\log n}$
(\secref{c05-contests}).

\solstep{Key Observations}
\begin{observation}
$x$'s parent is one of exactly two already-present keys: the largest key $<x$ (predecessor) or the
smallest key $>x$ (successor). Any other key is separated from $x$ by one of these two and could not
be the leaf's direct parent.
\end{observation}
\begin{observation}
Among the predecessor $p$ and successor $s$, the parent is whichever was inserted \textbf{later}.
The earlier one had the later already on the descent path, so $x$ attaches below the later one.
Equivalently: at every node, the child on the side $x$ belongs to was decided last among the two
neighbours.
\end{observation}
\begin{observation}
Both neighbours and their insertion times are available from a sorted \code{set} plus a
\code{map} key$\to$time in $\bigO{\log n}$ -- \emph{no tree is built}, so the degenerate shape never
costs us. This is the escape from the $\bigO{n^2}$ simulation.
\end{observation}

\solstep{Think It Through With Me}
\begin{thinkbox}
\deadend{Actually construct the BST with \secref{c05-insert-delete} insertions and print parents.}
{Correct but each insertion is an $\bigO{h}$ descent, and $h$ on sorted input is $n$, giving
$\bigO{n^2}$ -- the chain $C$ made manifest. For $n=10^{5}$ that is $\sim 10^{10}$ operations, a
\tle. The tree's \emph{bad shape} is exactly what we are trying to avoid; here we can avoid the tree.}
\deadend{Sort the array and answer neighbours by position.}
{A plain sort gives global neighbours but throws away \emph{insertion time}, and the parent depends on
which neighbour arrived \emph{later}. Sort order alone picks the wrong parent about half the time --
you need a structure that grows online, one key at a time.}
\aha{Maintain the present keys in a \code{std::set} (sorted, $\bigO{\log n}$ neighbours) and remember
each key's insertion index. $x$'s parent $=$ the predecessor or successor whose index is \emph{larger}.
One \code{lower\_bound} plus a step-back per key, $\bigO{n\log n}$ total, no tree, no degeneration.}
\end{thinkbox}

\solstep{Solution}
Iterate $i=1..n$ in insertion order. Before handling $a_i$, the \code{set} $s$ holds exactly the keys
already inserted (smaller than $i$ in time) and \code{when} maps each to its index. Compute
\code{r = s.lower\_bound(x)}. Candidate parent is \code{*r} (successor) when \code{r} is not
\code{s.end()}, and \code{*prev(r)} (predecessor) when \code{r} is not \code{s.begin()}; pick the
candidate with the larger \code{when}, which
is the later-inserted and hence the true parent. For $i\ge 2$ print that parent. Then insert $x$ into
$s$ and record \code{when[x]=i}. The root ($i=1$) prints nothing.

\solstep{Proof of Correctness}
\begin{proofbox}[The later-inserted neighbour is the parent]
Consider the moment $x$ is inserted and let $p<x<q$ be its predecessor and successor among the
already-present keys (one may be absent at an end). In the final BST every key of the tree lies in
exactly one of the open intervals determined by the sorted keys, and $x$'s search path is governed by
these: descending, the path is steered by $p$ and $q$ until it reaches whichever of them was inserted
last. Suppose $q$ was inserted after $p$ (the other case is symmetric). At the time $q$ arrived, $p$
was already present with $p<q$, so $p$ is in $q$'s left subtree; because no key lies strictly between
$p$ and $q$ (they are adjacent in the current sorted set), nothing between them can separate $x$ from
$q$, and $x$ (with $p<x<q$) descends at $q$ to the left child slot that is still empty -- hence $x$
becomes $q$'s left child, so $q$ is the parent. Conversely the earlier-inserted $p$ already had $q$
on the descent path below it, so it cannot be $x$'s direct parent. Therefore the parent is precisely
the neighbour with the greater insertion index, which the algorithm reports. The root ($i=1$) has no
neighbours and correctly prints nothing. \end{proofbox}

\solstep{Implementation}
\begin{lstlisting}[caption={Tree Construction -- set-based parents without building the tree (Codeforces 675D)}, label={lst:c06-treecons-soln}]
#include <bits/stdc++.h>
using namespace std;
int main(){
    ios::sync_with_stdio(false); cin.tie(nullptr);
    int n; cin >> n;
    vector<int> a(n);
    for (int i = 0; i < n; ++i) cin >> a[i];

    set<int> s;             // keys already inserted, kept sorted
    map<int,int> when;      // key -> insertion index (its "time")
    for (int i = 0; i < n; ++i) {
        int x = a[i];
        int parent = -1;
        auto r = s.lower_bound(x);            // first present key > x (distinct => never == x)
        // successor candidate *r, predecessor candidate *prev(r): pick the LATER-inserted one
        if (r != s.end())           parent = *r;
        if (r != s.begin()) {
            int p = *prev(r);                 // greatest key < x
            if (parent == -1 || when[p] > when[parent]) parent = p;
        }
        if (i > 0) cout << parent << '\n';    // root (i==0) has no parent
        s.insert(x);
        when[x] = i;
    }
    return 0;
}
\end{lstlisting}

\solstep{Trace}
\begin{tracebox}[Official-style sample: insert order $2, 1, 3, 5, 4$]
$i=0$, $x=2$: empty set, root, print nothing. $s=\{2\}$, \code{when[2]=0}.\\
$i=1$, $x=1$: \code{lower\_bound(1)}=iter to $2$ (successor). No predecessor ($r$ is begin). Parent
$=2$. Print $2$. $s=\{1,2\}$, \code{when[1]=1}.\\
$i=2$, $x=3$: \code{lower\_bound(3)}=end. Predecessor $=2$ (time $0$). Parent $=2$. Print $2$.
$s=\{1,2,3\}$, \code{when[3]=2}.\\
$i=3$, $x=5$: lb=end, pred $=3$ (time $2$). Parent $=3$. Print $3$. \code{when[5]=3}.\\
$i=4$, $x=4$: successor $=5$ (time $3$), predecessor $=3$ (time $2$). Later is $5$. Parent $=5$.
Print $5$. Output $2,2,3,5$ -- children $1,3,5,4$'s parents, matching the built tree. \checkmark
\end{tracebox}

\solstep{Complexity}
\begin{complexitybox}
Time $\bigO{n\log n}$: one \code{lower\_bound}, one \code{prev}, one \code{insert} per key, each
$\bigO{\log n}$ on a \code{std::set} of $\le n$ nodes (red-black -- \secref{c06-redblack} proves the
guarantee). Memory $\bigO{n}$ for the set and the time map. Crucially this is
$\bigO{n\log n}$ \emph{for every input order}, because we never let the tree's height enter the
cost -- the exact contrast with the $\bigO{n^2}$ simulation the section opened on. Verdict: \ac.
\end{complexitybox}

\solstep{Pitfalls and Variants}
\begin{warningbox}
Three ways to fall: (i) building the tree and descending (Trap~3) -- $\bigO{n^2}$ on sorted input.
(ii) forgetting distinctness -- here keys are distinct, so \code{lower\_bound} returns the first key
$\ge x$, which for a fresh $x$ is really the successor; with duplicates you must fix an
$\le$/$>$ convention. (iii) picking the \emph{earlier}-inserted neighbour -- inverted; it is the
\emph{later} one that is deeper and becomes the parent. \emph{Variants:} GNU \code{pb\_ds} or a hand
treap would also work but are overkill -- \code{set} already \emph{is} the balanced BST
(\secref{c06-redblack}). The same ``parent is the later neighbour'' idea reconstructs the whole tree,
not just the parent list.
\end{warningbox}

\begin{insightbox}[Key takeaways]
\begin{enumerate}[nosep, leftmargin=1.6em]
  \item \textbf{A rotation is the only legal move, and it is $\bigO{1}$ and inorder-preserving}
        (\thmref{c06-rotation}) -- every balanced tree is a \emph{policy} over rotations, never a new
        mechanic.
  \item \textbf{The shape is the enemy, not the data.} Same keys, same sorted line; a chain and a
        balanced tree differ only geometrically, and geometry decides $\bigO{h}$ vs $\bigO{n}$
        (\defref{c06-heightcost}).
  \item \textbf{You can often dodge the tree entirely.} When only sorted \emph{contents} and
        \emph{arrival times} matter, \code{std::set} gives $\bigO{\log n}$ without ever paying the
        degenerate height (\probref{c06-treecons}).
\end{enumerate}
\end{insightbox}

\begin{hardeningbox}
\begin{itemize}[nosep, leftmargin=1.4em]
  \item \textbf{Report depth instead of parent} -- with the same neighbours, a node's depth is one
        more than its (later-inserted) parent's; precompute times to answer both in one pass.
  \item \textbf{Reconstruct the whole tree} -- connect each key to its chosen parent and you have an
        $\bigO{n\log n}$ build of $C$, no descent.
  \item \textbf{Amortise later} -- splay (\secref{c06-splay}) rebuilds a BST from an order with the
        rotations themselves, giving $\bigO{n\log n}$ amortised \emph{while} physically holding a tree.
\end{itemize}
\end{hardeningbox}

\subsection{Practice Problems}

\begin{stuckbox}
\begin{itemize}[nosep, leftmargin=1.4em]
  \item Did you pick the \emph{later}-inserted neighbour, not the earlier, as the parent?
        (\probref{c06-treecons}.)
  \item Did you update heights \emph{bottom-first} after a rotation, so balance factors stay honest?
        (Trap~2, \coderef{c06-rotcode}.)
  \item For the rebuild practice: did you read the inorder (a sorted array) and then \emph{halve} it,
        rather than fiddling with rotations one by one?
\end{itemize}
\end{stuckbox}

\begin{practiceset}
  \practice{Balance a Binary Search Tree}{LeetCode 1382}
    {https://leetcode.com/problems/balance-a-binary-search-tree/}{LC Medium L3}
    {Given a possibly skewed BST, return \emph{any} height-balanced BST with the same values.
    The rotation-by-rotation route is fiddly; the clean route uses \thmref{c05-inorder}: dump the
    inorder into a sorted array, then build a balanced tree by recursively taking the middle element
    as the root (the same halving that makes a balanced tree from a sorted list). $\bigO{n}$ time and
    space, no rotations, and the sorted order is preserved by construction -- the extreme case of
    ``the shape is a choice'' (\thmref{c06-rotation}). Contrast with \probref{c06-treecons}, where the
    order was \emph{forced} and only the query was ours to answer.}
\end{practiceset}

\begin{recallbox}[10-second refresher -- 8.1 Why Balance? The Rotation]
Every BST operation is $\bigO{h}$ (\defref{c06-heightcost}), so keeping $h=\bigO{\log n}$ is the
whole game. The \term{rotation} re-links three pointers in $\bigO{1}$ and provably leaves the sorted
inorder unchanged (\thmref{c06-rotation}, \coderef{c06-rotcode}), so the shape is free to repair while
the data stays put. A same-side lean (LL, RR) needs one rotation; a zig-zag (LR, RL) needs two
(\defref{c06-rotcases}). And often the tree is not needed at all: CF 675D finds each parent as the
\emph{later-inserted} set neighbour in $\bigO{n\log n}$, sidestepping the degenerate height
(\probref{c06-treecons}).
\selftest{
  \item Draw the left rotation about $x$ and write the inorder of both pictures; why are they equal?
  \item The chain $1,2,3,4$ (all right children): how many rotations to balance, and which side each?
  \item Why is $Z$ (\figref{c06-zigzag}) not fixable by a \emph{single} rotation?
  \item In CF 675D, why is the \emph{earlier}-inserted neighbour provably not the parent?
  \item State the balance-height goal in one line: $h \le$ what, and why does it matter to search?
}
\end{recallbox}
